\section{Discussion}
\label{sec:discussion}

\subsection{What Makes a Property Easy to Order?}

Our results suggest that two factors determine ordering difficulty: (1) whether the items form a \emph{conventional sequence} that appears as an ordered list in training data, and (2) whether the ordering requires \emph{precise numeric discrimination} between items with similar values.

Factor (1) explains why element atomic numbers ($\tau = 1.0$) are easier than element atomic weights ($\tau = 0.970$): the periodic table is a fixed, universally-taught sequence, while atomic weights are independent numeric values that must be compared.
Factor (2) explains why country population ($\tau = 0.941$) is easier than country area ($\tau = 0.822$): population rankings are frequently discussed in media and education, while area comparisons between similarly-sized countries (\eg Australia vs.\ Brazil) require precise recall.

This framing aligns with \citet{zhao2025access}, who observe that pointwise ordering succeeds when facts reside in parametric memory.
We refine this observation: even when facts \emph{are} in parametric memory, the \emph{form} of storage matters.
Sequential knowledge (A comes before B in a known list) supports ordering more reliably than independent numeric knowledge (A has value $x$, B has value $y$, is $x > y$?).

\subsection{Error Patterns}

The most common errors are \emph{adjacent swaps} rather than gross misorderings.
For building height, \gptlarge swaps Big Ben (96\,m) and the Statue of Liberty (93\,m)---items separated by only 3 meters.
For country area, Australia and Brazil (similar sizes) are frequently transposed.
This pattern is consistent with LLMs storing \emph{approximate} magnitude representations that support coarse ordering but lose precision when values are close.

\gptsmall exhibits a qualitatively different failure mode on its worst tasks.
River length ($\tau = 0.274$) is near random ordering, suggesting that the smaller model has essentially no usable knowledge of relative river lengths.
This indicates that smaller models do not simply have noisier versions of the same knowledge---for some properties, the knowledge is absent entirely.

\subsection{Implications}

\para{For LLM-based data pipelines.}
Our hierarchy provides actionable guidance: trust LLMs for ordering by well-known categorical or sequential properties (alphabetical, chronological, famous rankings), but use external databases or APIs for ordering by precise numeric facts (geographic measurements, specific dimensions).

\para{For understanding LLM knowledge.}
Ordering tasks provide a clean, scalable probe of relational knowledge.
Unlike cloze-style probes that test whether an LLM can recall a single fact, ordering tests whether the model can \emph{compare} facts---a stronger requirement that reveals the precision of stored knowledge.
The gap between sequence-based and numeric ordering illuminates the boundary between ``knowing'' (recognizing a fact in context) and ``precisely recalling'' (producing a specific value that can be compared).

\para{For evaluation design.}
The \benchname hierarchy can guide benchmark construction.
Tasks that test sequence recall are easy and should not be used alone to assess factual knowledge.
Tasks that test precise numeric ordering are more discriminating and better suited for evaluating knowledge depth.

\subsection{Limitations}
\label{sec:limitations}

Our study has several limitations that temper the strength of our conclusions.

\para{Small task set.}
We test 22 tasks with 6--10 items each.
Larger lists would likely show more degradation, and more tasks per category would increase statistical power.
With only 2--7 tasks per category, between-category statistical tests lack power (Kruskal--Wallis $p = 0.10$ for \gptlarge).

\para{Two models.}
We evaluate only \gptlarge and \gptsmall.
Testing additional model families (Claude, Gemini, open-weight models like Llama) would strengthen generalizability claims.

\para{Ground-truth ambiguity.}
Some properties have measurement-dependent ground truths (\eg ``typical'' animal lifespan varies by source).
We use standard reference values, but alternative sources might yield different orderings.

\para{English only.}
All tasks and prompts are in English.
Cross-lingual ordering might reveal different accessibility patterns, especially for culture-specific knowledge.

\para{Temperature and prompting.}
We use temperature 0 and a minimal prompt.
Higher temperatures might reveal different variance patterns, and chain-of-thought prompting might improve performance on hard tasks---testing this would clarify whether the difficulty is in \emph{accessing} knowledge or in \emph{comparing} accessed values.

\para{Category boundaries.}
The distinction between ``well-known'' and ``specific'' facts is somewhat subjective.
Our results themselves suggest a better taxonomy: sequence-based versus independent-value properties.
